\begin{claim}{D.2.61}{Proved}{A secondary operation kills the remaining Moore obstruction}
Put $L=Y\langle7\rangle$, $E=\Omega^2L$, and let
$\widetilde\alpha:S^7\to L$ be the lift of $\alpha=u_4\nu_4$.
The bottom image of $h_0h_3$ in
$\operatorname{Ext}_{\mathcal A}^{2,16}(H^*L,\FF_2)$ is
a nonzero permanent cycle.  In particular
\[
                 2\widetilde\alpha\sigma\ne0
                       \quad\hbox{in }\pi_{14}^sL.
                 \tag{D.2.61a}
\]
The remaining stable Atiyah--Hirzebruch differential of
D.2.54 is zero on its whole source:
\[
 d_8:E^8_{15,0}\longrightarrow E^8_{7,7}=\ZZ/4
       \quad\hbox{is zero},\qquad E^\infty_{7,7}=\ZZ/4.
       \tag{D.2.61b}
\]
Consequently the actual Moore composite and the actual
sphere composite have the values
\[
                h_Ml_4=0,\qquad
                q=\widetilde\alpha\sigma'
                      \text{ has exact order two}.
                \tag{D.2.61c}
\]
The integer $n$ in D.2.53 is even, and the detecting
integral class of D.2.54 is a stable Hurewicz class.
These statements determine the specified Moore obstruction.
They do not construct the full fourth-suspension extension
or evaluate the nullification in D.2.4, which remains open.
\end{claim}
\begin{proof}
All unmarked cohomology in this proof has coefficients
$\FF_2$.  Write $a=a_7\in H^7L$ for the bottom class.
D.2.57 and instability give
\[
       \operatorname{Sq}^1a=\operatorname{Sq}^6a
                         =\operatorname{Sq}^8a=0.
       \tag{D.2.61d}
\]

\medskip\noindent
\textbf{The relation and a complete four-term functional.}
The Adem relations give
\[
 R=\operatorname{Sq}^1\operatorname{Sq}^8
   +\operatorname{Sq}^2\operatorname{Sq}^1\operatorname{Sq}^6
   +\operatorname{Sq}^8\operatorname{Sq}^1=0.
   \tag{D.2.61e}
\]
Indeed the first term is $\operatorname{Sq}^9$, and the
middle term is
$\operatorname{Sq}^2\operatorname{Sq}^7
 =\operatorname{Sq}^9+\operatorname{Sq}^8\operatorname{Sq}^1$.
In the normalized bar complex dual to the left-comodule
cobar convention of D.2.57, consider
\[
 r=[\operatorname{Sq}^1|\operatorname{Sq}^8]a
  +[\operatorname{Sq}^2\operatorname{Sq}^1|
                                  \operatorname{Sq}^6]a
  +[\operatorname{Sq}^8|\operatorname{Sq}^1]a.
  \tag{D.2.61f}
\]
The bar differential is
$d([b|c]m)=[bc]m+[b](cm)$.
Thus \textup{(D.2.61d)--(D.2.61e)} make $r$ a cycle.
Its pairing defines a functional $\lambda_R$ on
$\operatorname{Ext}_{\mathcal A}^{2,16}(H^*L,\FF_2)$.
In Milnor coordinates its entire support is
\[
 [\xi_1|\xi_1^8]A,\quad
 [\xi_1^3|\xi_1^6]A,\quad
 [\xi_2|\xi_1^6]A,\quad
 [\xi_1^8|\xi_1]A,
 \tag{D.2.61g}
\]
where $A$ is dual to $a$.  In particular
\[
              \lambda_R([\xi_1|\xi_1^8]A)=1.
              \tag{D.2.61h}
\]
This is a symbolic certificate on the full module.
Independently, the authored script
\texttt{code/EH5\_secondary\_relation\_certificate.py}
checks these four terms against all 166 incoming columns
of the actual degree-fifteen Milnor input and records
every supported term in
\texttt{code/outputs/EH5\_secondary\_relation\_certificate.json}.
No secondary topological assertion is inferred from that
finite algebra check alone.

\medskip\noindent
\textbf{An actual universal choice of the secondary operation.}
Let $H=H\FF_2$ and let $\mathcal P$ be the fiber of the
stable map
\[
 (\operatorname{Sq}^1,\operatorname{Sq}^6,\operatorname{Sq}^8):
 \Sigma^7H\longrightarrow
           \Sigma^8H\oplus\Sigma^{13}H\oplus\Sigma^{15}H.
 \tag{D.2.61i}
\]
The row
$(\operatorname{Sq}^8,\operatorname{Sq}^2\operatorname{Sq}^1,
\operatorname{Sq}^1)$ lands in $\Sigma^{16}H$.
A fixed nullhomotopy of its composite, supplied by $R$,
defines a secondary map $\mathcal P\to\Sigma^{15}H$.
Its restriction to the fiber of
$\mathcal P\to\Sigma^7H$ is the loop of that row.
This is the universal construction of secondary operations;
the precise fiber restriction, existence, and uniqueness
up to a primary operation are given by
\cite[Section~3.6, especially Theorems~3.6.1--3.6.2 and
Lemma~3.6.3, pp.~68--74]{EH5Adams1960Secondary}.

Use the actual space
\[
 \mathcal F=\operatorname{hofib}
  \bigl(\operatorname{Sq}^6\rho:
       K(\ZZ/4,7)\longrightarrow K(\FF_2,13)\bigr).
 \tag{D.2.61j}
\]
Here $\rho$ denotes reduction modulo two.
On the base $K(\ZZ/4,7)$ choose based nullhomotopies
of $\operatorname{Sq}^1\rho$ and
$\operatorname{Sq}^8\rho$.
They exist respectively from the mod-four lift and
instability.  Together with the tautological fiber path
for $\operatorname{Sq}^6\rho$, they give a map
$\mathcal F\to\Omega^\infty\mathcal P$ and hence a
particular class $\phi\in H^{15}\mathcal F$.
On its fiber $K(\FF_2,12)$ the two base-chosen
nullhomotopies are constant.  Therefore
\[
             \phi|_{K(\FF_2,12)}
                    =\operatorname{Sq}^2\operatorname{Sq}^1
                                                \iota_{12}.
             \tag{D.2.61k}
\]
This fixes an actual representative; it does not replace
that restriction by a coset allowing an arbitrary
$\operatorname{Sq}^3\iota_{12}$ term.

Choose a mod-four lift $\widehat a\in H^7(L;\ZZ/4)$
of $a$.  Such a lift exists by
$\operatorname{Sq}^1a=0$.  Since
$\operatorname{Sq}^6a=0$, it lifts to an actual map
$f:L\to\mathcal F$.
The class $f^*\phi$ is an admissible value of the
degree-eight operation $\Phi_R(a)$.
Its indeterminacy is
\[
 \operatorname{Sq}^8H^7L+
 \operatorname{Sq}^2\operatorname{Sq}^1H^{12}L+
 \operatorname{Sq}^1H^{14}L.
 \tag{D.2.61l}
\]
Every class annihilated by positive homology Steenrod
operations annihilates this entire indeterminacy.

Double looping \textup{(D.2.61j)} gives an actual splitting
\[
 \Omega^2\mathcal F\simeq
            K(\ZZ/4,5)\times K(\FF_2,10)
 \tag{D.2.61m}
\]
over $K(\ZZ/4,5)$: the defining map is now
$\operatorname{Sq}^6$ on a degree-five class and is null.
Choose one such splitting.  There are no mixed positive
products in degree thirteen, since $5+10>13$.
Writing $z_{10}$ for the second factor's fundamental
class, \textup{(D.2.61k)} implies
\[
 \sigma^2\phi=\operatorname{Sq}^2\operatorname{Sq}^1z_{10}
                       +\eta,
       \qquad \eta\in H^{13}K(\ZZ/4,5).
 \tag{D.2.61n}
\]
Serre's computation
\cite[Section~2, no.~11, Theorem~4 and its proof,
pp.~204--205]{I_Serre1953}
says every polynomial generator of this last algebra
is a Steenrod operation on the fundamental reduction
$x_5$ or on $y_6=\rho\beta_4\widehat x_5$.
In degree thirteen, $\eta$ is therefore a sum of
positive-primary images of $x_5,y_6$ and decomposables.
The higher coefficient connecting map here is the one
for $0\to\FF_2\to\ZZ/8\to\ZZ/4\to0$;
it is $\rho\beta_4$ and commutes with suspension.
Under $\Omega^2f$, both $x_5$ and $y_6$ are double
cohomology suspensions of the actual classes
$a$ and $\rho\beta_4\widehat a$ on $L$.

\medskip\noindent
\textbf{Evaluation on every possible Adams source.}
Let $x\in H_{15}(L;\FF_2)$ be annihilated by every
positive homology Steenrod operation.
D.2.58 proves that $\sigma^2:H_{13}E\to H_{15}L$
is onto.  We can choose a primitive preimage $z$.
Indeed $E$ is four-connected, so its degree-thirteen
reduced diagonal has only bidegrees $(5,8),(6,7)$
and their reverses.  All factors in those degrees are
primitive.  The nine products
$uH_8E$ and $H_6E\,H_7E$ have independent symmetric
diagonals and realize all these tensors.  Subtracting
a suitable sum of them makes a preimage primitive;
these products have zero double homology suspension.

Naturality makes $\operatorname{Sq}^2_*z$ primitive
and places it in the complete D.2.56 kernel
\[
 \ker(\sigma^2:H_{11}E\to H_{13}L)
                    =\langle uD,uE_6,\xi u\rangle.
 \tag{D.2.61o}
\]
The first two products have independent nonzero
reduced diagonals, so only $\xi u$ could remain.
Its coefficient is zero already in the ambient free
$E_2$ homology of $T=\Omega^2\Sigma^2B$.
The only degree-six restricted-Lie inputs for a
degree-thirteen restricted generator are $f,x,\{a,b\}$;
their Bocksteins are $0,d+e,0$.
Thus the top Nishida formula cannot produce the outer
restricted generator $\xi u$, where $u=\xi a$.
Ordinary labels, Lie words and products cannot produce
that outer restricted generator either.
These are the free-$E_2$ basis and top-square formulas
already used in D.2.56--D.2.58
\cite[Chapter~III, Theorem~1.3(3), pp.~217--218]{I_CohenLadaMay1976}.
The established injection $H_{11}E\to H_{11}T$ gives
\[
                         \operatorname{Sq}^2_*z=0.
                         \tag{D.2.61p}
\]
This argument uses no new degree-thirteen cover-edge
calculation and does not assert $\operatorname{Sq}^1_*z=0$.

Pair \textup{(D.2.61n)} with $(\Omega^2f)_*z$.
The decomposable part of $\eta$ pairs zero by primitivity.
Every positive-primary part is a double suspension of
a positive operation on one of the two stated classes
of $L$, so it pairs zero with the primary-annihilated
image $x$.  The fiber term pairs as
$\operatorname{Sq}^1_*\operatorname{Sq}^2_*z=0$.
Consequently
\[
             \langle\Phi_R(a),x\rangle=0
            \quad\text{for every such }x.
             \tag{D.2.61q}
\]
The controlled universal representative is essential:
arbitrary secondary indeterminacy after looping could
include $\operatorname{Sq}^1H^{12}E$, which need not
pair zero with the chosen $z$.

\medskip\noindent
\textbf{The precise secondary-operation/Adams comparison.}
We verify the needed comparison rather than identifying
operations merely by their degree.  Put $M=H^*L$.
Use the actual GEM Adams coresolution of
$X=\Sigma^\infty L$, identified with the nonminimal
free bar resolution
\[
 P_0=\mathcal A\otimes M,\qquad
 P_1=\mathcal A\otimes\overline{\mathcal A}\otimes M,
 \qquad
 P_2=\mathcal A\otimes\overline{\mathcal A}^{\otimes2}
                                                   \otimes M.
 \tag{D.2.61r}
\]
Its coherent zero tracks come from the actual Adams
fiber tower, as in
\cite[Theorem~6.9 and Section~7]{EH5BauesJibladze2004};
they are not assigned solely from a primary chain complex.
Here the indicated free resolution can be realized explicitly.
Choose $g_0:X\to A_0$ to induce its augmentation
$P_0\to M$.  Its cofiber $C_0$ has cohomology
$\ker(P_0\to M)$, because the augmentation is onto
in every degree.  Realize the free surjection from $P_1$
onto this kernel by $C_0\to A_1$, and repeat with its
cofiber and the free surjection from $P_2$.
The resulting cofiber triangles give the required zero
tracks and their coherence.  There are only finitely many
generators in each bounded degree, so the connective GEM
products and these cohomology identifications are legitimate.
Write the corresponding GEM maps as
$X\xrightarrow{g_0}A_0\xrightarrow{v_0}A_1
\xrightarrow{v_1}A_2$, with null tracks
$\eta_0$ on $v_0g_0$ and $\eta_1$ on $v_1v_0$.
Coherence says $\eta_1g_0=v_1\eta_0$.

In the free resolution,
\[
 d_P([b|c]m)=b[c]m+[bc]m+[b](cm).
 \tag{D.2.61s}
\]
For the lift of $r$ in \textup{(D.2.61f)}, the last
two sums vanish by \textup{(D.2.61d)--(D.2.61e)}.
Its coordinate of $v_1$ is therefore exactly the row
$(\operatorname{Sq}^8,
\operatorname{Sq}^2\operatorname{Sq}^1,
\operatorname{Sq}^1)$ on the three coordinates of $A_1$
in degrees $8,13,15$.
The expression $r$ is a nonzero scalar combination of
three free $P_2$ generators, so an ordinary change of
free basis makes it a single coordinate as used here.
Let $H_R$ be the chosen relation nullhomotopy on $A_0$
factoring through its bottom $\Sigma^7H$ coordinate.
The difference of tracks
$\eta_{1,r}-H_R$ is a stable map
$q_0:A_0\to\Sigma^{15}H$.
Coherence gives
\[
                       g_0^*q_0=\Phi_R(a)
                         \quad\text{modulo primary indeterminacy}.
                       \tag{D.2.61t}
\]

A class in $\operatorname{Ext}^{0,15}_{\mathcal A}(M,\FF_2)$
is exactly a primary-annihilated $x$ as above.
Represent it by $c:S^{15}\to A_0$, chosen with nonzero
components only in degree-fifteen GEM coordinates, and
choose a null track $\gamma$ for $v_0c$.
The actual secondary differential is represented by
the glued track $\eta_1c+v_1\gamma$.
This is the construction in
\cite[equation~(4.1), Theorem~4.2, and Theorem~7.3 with
its proof, pp.~13--15 and 19--21]{EH5BauesJibladze2004},
which identifies it with the Adams $d_2$.
On the $r$ coordinate, the relevant components of
$\gamma$ map to
$\Omega\Sigma^8H$, $\Omega\Sigma^{13}H$ and
$\Omega\Sigma^{15}H$; their groups of maps from
$S^{15}$ are all zero.  The track $H_Rc$ is constant,
since the chosen $c$ has zero bottom $\Sigma^7H$
component.  Thus only $q_0c$ remains, giving the exact
pairing formula
\[
       \lambda_R(d_2x)
          =\langle g_0^*q_0,x\rangle
          =\langle\Phi_R(a),x\rangle.
       \tag{D.2.61u}
\]
Changing the three primary nullhomotopies adds precisely
\textup{(D.2.61l)}.  Changing $H_R$ adds a degree-eight
primary operation on $a$, by the same universal
construction or Adams's uniqueness theorem.
Every such change pairs zero with $x$.
This accounts for both the source-dependent tracks
and the full indeterminacy in \textup{(D.2.61u)}.

\medskip\noindent
\textbf{Permanence and the actual conclusion.}
Equations \textup{(D.2.61h)}, \textup{(D.2.61q)} and
\textup{(D.2.61u)} show that no incoming Adams $d_2$
can hit the bottom $h_0h_3$ image.
There is no later incoming differential to filtration
two, and sphere naturality prevents outgoing
differentials, as in D.2.58.
The connective finite-type Adams convergence theorem
used in D.2.57 therefore detects a nonzero image of the
actual bottom $2\sigma$ in the two-completed stable
homotopy of $L$
\cite[Part~III, Theorem~15.1 and its construction,
pp.~370--377]{EH5Adams1974AHSS}.
A zero actual stable class would have zero completed
image, so \textup{(D.2.61a)} follows without requiring
a separate completion-isomorphism assertion.

The bottom subgroup $E^8_{7,7}$ is $\ZZ/4$ by D.2.57.
Its twice-generator is the actual class detected in
\textup{(D.2.61a)}.  Any nonzero subgroup of this
$\ZZ/4$ contains that element.  Thus a nonzero incoming
$d_8$ would kill the actual bottom stable class.
There is no smaller cell filtration, since $L$ is
six-connected; hence it cannot disappear into a lower
filtration term.  This proves $d_8=0$ on its whole
source.  No later incoming differential is possible
because the sphere spectrum is connective, and no
outgoing differential from the bottom filtration is
possible.  This proves \textup{(D.2.61b)}.

The detecting class $\omega$ of D.2.54 now survives.
It has no later possible outgoing differential and no
incoming differential in the $q=0$ row.  Finite-range
cellular convergence makes it the integral Hurewicz
image of a stable map.  Its $\theta$ pairing is one,
so the exact D.2.51 criterion gives $h_Ml_4=0$.
Equivalently, the exact D.2.54 equation gives
$2n\bar\alpha=0$ in $\ZZ/4$, so $n$ is even and the
D.2.54 source-changing criterion holds with zero change.
Finally D.2.53 gives $2q=2h=0$, whereas the stabilization
of $q=\widetilde\alpha\sigma'$ is the nonzero class
$2\widetilde\alpha\sigma$ of \textup{(D.2.61a)}.
Thus $q$ has exact order two.
None of these implications supplies an extension over
all the other cells of the fourth-suspension source;
the main D.2.4 test is unchanged.
\end{proof}
